\section{The canonical cyclotomic tower in every finite class}
\label{sec:canonical-all-stages}

Return to the actual cyclotomic source
\[
 H=\operatorname{Gal}(\Omega_p/F_\infty),\qquad
 M=H^{\mathrm{ab}},\qquad
 V=M_{\Lambda\text{-tor}}=D\oplus T,\qquad W=M/V.
\]
The arithmetic results already established give that \(H\) is free
pro-\(p\), \(V\) is finite free over \(\mathbb Z_p\), \(W\) is
torsion-free over \(\Lambda\), and \(W[1/p]\) is finite free over
\(\Lambda[1/p]\).  These use Ferrero--Washington and Iwasawa's
freeness theorem as recorded in \S\ref{sec:abelian}.
Fix the characteristic polynomial
\[
               a(Z)=\det(Z-\gamma\mid V).
\]
Lemma~\ref{lem:positive-eigenvalues} proves
\begin{equation}\label{eq:uniform-positive-exclusion}
                       a(u^m)\ne0\qquad(m\geq2).
\end{equation}

Here is a recursion involving only the source and its arithmetic
action.  Put \(\mathcal F_1=H\), \(\mathcal R_1=M\), and let \(\mathcal F_2\) be the inverse
image of \(V\).  For \(n\geq2\), put
\begin{equation}\label{eq:uniform-recursion}
 J_n=\overline{\prod_{\substack{i+j=n\\i,j\geq1}}[\mathcal F_i,\mathcal F_j]},
 \qquad
 \mathcal R_n=\mathcal F_n/J_{n+1},
 \qquad
 \mathcal F_{n+1}=\operatorname{preim}_{\mathcal F_n}
                  \bigl((\mathcal R_n)_{\Lambda\text{-tor}}\bigr).
\end{equation}
The proof below establishes at each step that these expressions
are defined and that the displayed torsion subgroup is closed.
Set
\[
 Q_n=H/\mathcal F_{n+1},\qquad E_n=\mathcal F_n/\mathcal F_{n+1}.
\]

We prove Theorem~\ref{thm:canonical-all-stages}. In the course of the proof, the map from \((\mathcal R_n)_{\Lambda\text{-tor}}\) to \(M\) will also be shown to have kernel killed by \(p^{h_n}\).

The exact source-intersection question remaining in the theorem is whether
\(\bigcap_n\mathcal F_n=1\).  The source countermodel below shows why the
finite-stage assertions alone do not imply that equality.

\subsection{A uniform source calculation}

We give the integral comparison needed in the induction.
The statement distinguishes the actual commutator surjection
from the rational Hall isomorphism.

Suppose the construction is known through \(Q_{n-1}\).  Then
\([\mathcal F_i,\mathcal F_j]\subseteq \mathcal F_{i+j}\) whenever the right side has
already been constructed.  This follows either from its
established target description or from the corresponding
central filtration.  Thus \(J_{n+1}\subseteq \mathcal F_n\), and \(\mathcal R_n\)
is central in \(H/J_{n+1}\).  There is an actual exact sequence
\begin{equation}\label{eq:uniform-central-sequence}
 0\longrightarrow \mathcal S_n:=J_n/J_{n+1}
 \longrightarrow \mathcal R_n
 \longrightarrow (\mathcal R_{n-1})_{\Lambda\text{-tor}}
 \longrightarrow0.
\end{equation}
Indeed \(\mathcal F_n/J_n\) is exactly the torsion subgroup used at the
preceding step.

\paragraph{Rational Hall coordinates with a common bound.}
Work in \(H/\gamma_{n+1}H\).  A compact
\(\mathbb Z_p\)-basis of \(M\), using a continuous \(\mathbb Z_p\)-module splitting \(M\simeq W\oplus V\), lifts to a convergent
free basis of \(H\), by the Frattini and projectivity argument
used above.  Its free nilpotent Hall coordinates identify
\[
       \log(H/\gamma_{n+1}H)[1/p]
\]
with the completed free rational Lie algebra on \(M[1/p]\),
truncated at ordinary length \(n\).
In fixed length, BCH, logarithm, and exponential use a finite
list of universal rational coefficients.  One integer clears
their denominators independently of the number of basis
coordinates.  Passing through the finite-tail quotients
therefore gives the same common bound on the compact completion.

Let
\[
 A_i=\operatorname{im}(\mathcal F_{i+1}\longrightarrow M)\subseteq V.
\]
Choose a rational basis of \(V[1/p]\) adapted to the finite flag
\[
 V[1/p]=A_1[1/p]\supseteq\cdots\supseteq A_{n-1}[1/p].
\]
Assign weight \(j\) to a complement of \(A_j[1/p]\) in
\(A_{j-1}[1/p]\), and weight \(n\) to the remaining space.
Assign weight one to \(W[1/p]\).
The finite matrices and their inverses have a common denominator.
After scaling, the selected finite basis vectors lift to the
indicated \(\mathcal F_j\).

The literal source maps are obtained as follows.  The image of
a weight-\(j\) lift in \(Q_{j-1}\) has zero abelianization.
Express it in the already established derived Hall coordinates
of weights at most \(j-1\), lift that expression to the free
source, and subtract it.  The corrected generator lies in the
rational kernel of the actual map to \(Q_{j-1}\), with unchanged
abelian image.  Correct the remaining weight-\(n\) vectors
against \(Q_{n-1}\) itself.  There are finitely many such
corrections.  The bounded Hall sections at the preceding
stages and the fixed-class BCH bound give one bound for them.

These substitutions have invertible bounded linear part and
derived corrections.  Their inverse is obtained successively
in ordinary degrees at most \(n\), again with a common bound.
The resulting map to \(Q_{n-1}\) is the literal weighted
truncation map.  To see this, its map on each already
established relative graded piece is the free Hall
isomorphism; a finite descending-filtration argument gives
an isomorphism on the truncation.  The same statement holds
for every earlier quotient.  Thus the rational image of
\(\mathcal F_i\) is exactly the ideal of assigned weights at least
\(i\), as an equality of kernels of the actual maps.

Every nonlinear monomial of assigned weight at least \(n+1\)
is a bracket of monomials whose weights dominate positive
\(i,j\) with \(i+j=n+1\); hence it belongs to \(J_{n+1}\).
The reverse inclusion follows from the assigned weights of
its defining brackets.  Consequently
\begin{equation}\label{eq:uniform-rational-Hall}
 \mathcal S_n[1/p]\simeq
 \operatorname{Lie}^{\mathrm{wt}}_n
      \bigl(W[1/p],C_2,\ldots,C_{n-1}\bigr),
\end{equation}
where the \(C_j\) are the finite rational generator spaces
already recovered in weight \(j\).
The only additional rational generators in \(\mathcal R_n\) are the
remaining abelian vectors assigned weight \(n\).

\paragraph{The integral central presentation and its torsion.}
Form
\[
 \mathcal P_n=
 \bigoplus_{\substack{i<j\\i+j=n}}
       E_i\widehat\otimes_{\mathbb Z_p} E_j
 \ \oplus\
 \begin{cases}
 \widehat{\bigwedge}^{\,2}E_{n/2},&n\text{ even},\\
 0,&n\text{ odd}.
 \end{cases}
\]
Commutators give a continuous surjection
\(\mathcal P_n\to \mathcal S_n\).  Changing a representative introduces
a commutator in \(J_{n+1}\); the resulting central brackets
are \(\mathbb Z_p\)-bilinear.  Quotient by all Jacobi images
for triples of positive weights summing to \(n\), using the
actual lower bracket maps.  Hall--Witt gives
\begin{equation}\label{eq:uniform-Hall-surjection}
       \mathcal C_n:=\mathcal P_n/\mathcal I_n
                         \twoheadrightarrow \mathcal S_n.
\end{equation}
Equation~\eqref{eq:uniform-rational-Hall} makes this a rational
isomorphism.  Its integral kernel may contain inherited
bounded \(p\)-torsion.

Choose bracket-compatible compact free weighted Hall lattices
\(E_{i,0}\subseteq E_i\), \(i<n\), with
\[
 p^{b_i}E_i\subseteq E_{i,0},\qquad b_1=0.
\]
They are generated by the actual compact \(W\)-lattice and
finitely many selected lifts of each recovered \(C_j\).
Define
\begin{equation}\label{eq:uniform-bounds}
 d_n=\max_{i+j=n}(b_i+b_j),\qquad
 j_n=\max_{i+j+k=n}(b_i+b_j+b_k),
\end{equation}
taking \(j_2=0\).  Then
\[
 p^{d_n}\mathcal P_n\subseteq\mathcal P_{n,0},\qquad
 p^{j_n}\mathcal I_n\subseteq\mathcal I_{n,0}.
\]
The free Hall quotient
\(\mathcal P_{n,0}/\mathcal I_{n,0}\) is
\(p\)-torsion-free.  If \(x\) represents torsion in
\(\mathcal C_n\), some \(p^kx\) lies in \(\mathcal I_n\).
The element \(p^{d_n}x\) belongs to \(\mathcal P_{n,0}\),
and a further power \(p^{k+j_n}\) puts it in
\(\mathcal I_{n,0}\).  Saturation of the free Hall relation
module gives
\begin{equation}\label{eq:uniform-S-torsion}
 p^{d_n}\mathcal C_n[p^\infty]=0,\qquad
 p^{d_n}\mathcal S_n[p^\infty]=0.
\end{equation}
The second assertion follows because the kernel of
\eqref{eq:uniform-Hall-surjection} is torsion.
Also \(p^{d_n}\mathcal S_n\) lies in the free Hall image.
The pure module \(\operatorname{Lie}_n(W)\) injects into
the actual \(\mathcal S_n\), since its map is rationally injective
and its source is \(p\)-torsion-free.

\subsection{Induction against the actual outer image}

The rational outer image supplied by Hain--Matsumoto and
Brown is the weight-completed free Lie image on the odd
letters \(s_m\), \(m\geq3\), as used in the preceding
sections \cite{HM,Brown}.  Use its own generator-length
filtration.  After the established normalization,
\(\gamma s_m=u^m s_m\).  Every finite weight piece of
its length-\(n\) quotient is a sum of positive-weight
eigenspaces.

The base \(n=1\) is the already proved moment detection:
the degree-one coefficient map has kernel \(V\), so
\[
 \mathcal F_2=\ker\rho_{\leq1},\qquad Q_1=W.
\]
The torsion formula has \(h_1=0\).

Assume the theorem through \(n-1\).
The source calculation gives
\eqref{eq:uniform-rational-Hall}.  Apply
Theorem~\ref{thm:weighted-companion-saturation} to the
finite alphabet already recovered through weight \(n-1\).
Its unrestricted degree-\(n\) coefficient map is injective,
and its total-divisor map is saturated.

These are the actual normalized target coefficients.
Indeed the action of \(\gamma\) on a finite recovered
generator and its companion gives
\eqref{eq:wcs-resolvent} on every ordered odd-weight tuple.
Every finite generator is a subquotient of \(V[1/p]\);
its companion polynomial therefore has roots among those
of \(a\).  Equation~\eqref{eq:uniform-positive-exclusion}
ensures that none of its denominators vanishes on the
positive-weight evaluation grid.  Clearing the common
denominator reduces vanishing to the bounded multivariable
moment separation already proved in the class-two
comparison.  It follows that
\begin{equation}\label{eq:uniform-S-injection}
 \mathcal S_n[1/p]\longrightarrow
   \operatorname{gr}^{\mathrm{length}}_n
                          \operatorname{Lie}\mathcal U
\end{equation}
is injective.  In particular \(\mathcal S_n[1/p]\) is
\(\Lambda[1/p]\)-torsion-free; its integral
\(\Lambda\)-torsion is precisely its bounded
\(p\)-torsion from \eqref{eq:uniform-S-torsion}.

Put \(h_n=h_{n-1}+d_n\).  The quotient in
\eqref{eq:uniform-central-sequence} is annihilated by
\(p^{h_{n-1}}a(\gamma)\).  This multiple of a torsion
element of \(\mathcal R_n\) consequently lies in the bounded
torsion of \(\mathcal S_n\).  Conversely, an element killed
by a nonzero element of \(\Lambda\) is torsion by
definition.  Hence
\begin{equation}\label{eq:uniform-closed-torsion}
 (\mathcal R_n)_{\Lambda\text{-tor}}
           =\ker\bigl(p^{h_n}a(\gamma):\mathcal R_n\to \mathcal R_n\bigr).
\end{equation}
This proves closedness before taking its quotient.
The action of \(H\) on \(\mathcal R_n\) is trivial, and the
arithmetic action preserves its \(\Lambda\)-torsion.
Thus \(\mathcal F_{n+1}\) is a closed arithmetically stable
normal subgroup.  The quotient
\(E_n=\mathcal R_n/(\mathcal R_n)_{\Lambda\text{-tor}}\) has no
\(\Lambda\)-torsion, including no \(p\)-torsion.

The map \((\mathcal R_n)_{\Lambda\text{-tor}}\to M\) has kernel
killed by \(p^{h_n}\).  An element of that kernel maps,
in \eqref{eq:uniform-central-sequence}, to the preceding
kernel; multiplication by \(p^{h_{n-1}}\) puts it in
\(\mathcal S_n\), and multiplication by \(p^{d_n}\) then kills it.
Its image is therefore a closed finite-rank submodule
\(A_n\subseteq V\), and the new rational generator space is
\[
                         (A_{n-1}/A_n)[1/p].
\]

It remains to verify the exact target kernel.
The actual length-\(n\) map is defined on \(\mathcal R_n\),
because every bracket defining \(J_{n+1}\) has target
length at least \(n+1\).  Its value on a torsion element
is killed by \(p^{h_n}a(\gamma)\).
On every finite geometric weight piece this operator
is invertible by \eqref{eq:uniform-positive-exclusion};
thus torsion maps to zero in the full target.

Conversely, if \(r\in \mathcal R_n\) has zero length-\(n\) image,
then \(p^{h_{n-1}}a(\gamma)r\) lies in \(\mathcal S_n\) and
still has zero image.  By
\eqref{eq:uniform-S-injection} it is \(p\)-torsion,
and \eqref{eq:uniform-S-torsion} kills it.  Thus
\(r\in(\mathcal R_n)_{\Lambda\text{-tor}}\).
We have proved
\[
                    \mathcal F_{n+1}=\ker\rho_{\leq n}.
\]
Compactness gives the asserted isomorphism with the
actual image.

For the next induction step, the rational \(E_n\) is
the primitive companion extension of
\(\mathcal S_n[1/p]\) by the finite space
\((A_{n-1}/A_n)[1/p]\).
Its torsion-freeness and the total saturation of \(\mathcal S_n\)
give independent companion residues.  In the associative
algebra, the Eulerian idempotent decomposes the degree
module as this actual primitive extension plus the
nonprimitive old summand.  Both are torsion-free.
Thus all hypotheses of the weighted companion theorem
are verified for the enlarged alphabet.

We also close the integral lattice induction.
The quotient \(E_n/\operatorname{im}\mathcal S_n\) has bounded
\(p\)-torsion and finite free quotient.  Its bounded
torsion comes from the preceding kernel, already killed
by \(p^{h_{n-1}}\), and from the finite torsion in the
quotient of two finite-rank submodules of \(V\).
Choose \(p^{c_n}\) killing it and lift a basis of its
finite free quotient.  These finitely many lifts,
together with the old free Hall image, define
\(E_{n,0}\).  Since \(p^{d_n}\mathcal S_n\) lies in that Hall
image,
\[
                 p^{b_n}E_n\subseteq E_{n,0},
                 \qquad b_n=d_n+c_n.
\]
This proves the bounds required at the next stage.

Finally, \(Q_n\) is torsion-free because its finite
central filtration has the \(p\)-torsion-free factors
\(E_1,\ldots,E_n\).  It has class at most \(n\).
The integral pure Hall module
\(\operatorname{Lie}_n(W)\subseteq \mathcal S_n\) survives
in \(E_n\): its intersection with the torsion kernel
is zero.  This module is nonzero, since the positive
\(\Lambda\)-rank of \(W\) supplies two independent
\(\mathbb Z_p\)-coordinates and their length-\(n\)
Hall commutator.  Therefore the class is exactly \(n\).

\subsection{The full outer kernel and the inverse limit}

We justify the equality of the rational and pro-\(p\)
outer kernels used in the limit.  For the geometric
free pro-\(p\) group \(\Pi=\langle x,y\rangle\),
the tangential normalization is
\[
 x\longmapsto x^{\chi(\sigma)},\qquad
 y\longmapsto f_\sigma^{-1}y^{\chi(\sigma)}f_\sigma,
 \qquad f_\sigma\in[\Pi,\Pi].
\]
This normalized action on the cyclotomic kernel of \(G_{\mathbb Q}\) is given in
\cite[\S2, (2.1) and the following paragraph, author
p.~7]{NakamuraSchneps2000}. Here \(\chi=1\). The kernel comparison below also proves its descent to \(H\): an element trivial in the pro-\(p\) outer action has identity normalized action, so the kernel of the maximal extension unramified outside \(p\) acts trivially.

The map from \(\Pi\) to its rational prounipotent
completion is injective.  One way to verify this
directly is the Magnus map \(x\mapsto1+X\),
\(y\mapsto1+Y\).  Its reduction modulo \(p\)
detects every finite \(p\)-quotient: substitution
into the augmentation algebra of that quotient
is defined because its augmentation ideal is
nilpotent.  Hence its kernel in the free pro-\(p\)
group is trivial.  Inclusion of integral
coefficients into rational coefficients, followed
by the invertible formal exponential change of
variables, gives the stated rational injection.

Suppose a normalized cyclotomic action
is inner in the rational completion.  A conjugating
element fixing \(\exp X\) is \(\exp(aX)\):
the centralizer of \(X\) in the completed free
associative algebra is \(\mathbb Q_p[[X]]\),
by the following coefficient argument. In a homogeneous series commuting with \(X\), the coefficient of a word \(X^rYv\) equals that of \(YvX^r\) by successively moving the initial \(X\) to the end. Comparing \(XA=AX\) at the word beginning in \(Y\) and ending in \(X\) makes this coefficient zero. Thus every word containing \(Y\) has zero coefficient. Moreover primitivity in \(\mathbb Q_p[[X]]\) is the identity \(f(X_1+X_2)=f(X_1)+f(X_2)\); comparison of powers in characteristic zero leaves only its linear term. The primitive elements are precisely
\(\mathbb Q_pX\).
Use the convention that the inner action is \(z\mapsto g^{-1}zg\). Then \(g=\exp(aX)\), and the formula on \(y\) says that \(f_\sigma\exp(-aX)\) centralizes \(\exp Y\). Hence \(f_\sigma=\exp(bY)\exp(aX)\).
The abelian image of \(f_\sigma\) is zero;
abelianization forces \(a=b=0\).
The Magnus injection forces \(f_\sigma=1\)
in \(\Pi\).  Thus the rationally trivial
outer action is already trivial as a pro-\(p\)
outer action.  The reverse implication is
functorial, proving equality of the kernels.

The same conclusion holds when rational outer triviality is formulated at every finite geometric quotient. In the free Lie algebra truncated at degree \(N\), the centralizer of \(X\) is \(\mathbb Q_pX\) plus the top-degree center, by the same coefficient argument in degrees below \(N\). A conjugator therefore has logarithm \(aX\) modulo that center. The condition \(f_\sigma\in[\Pi,\Pi]\) makes the normalized action have no degree-two term on \(Y\), whereas this conjugator gives \(-a[X,Y]\). For \(N\geq2\), this forces \(a=0\), and the inner action is the identity. Passing through the quotients proves the assertion without choosing compatible conjugators.

The generator-length filtration of the
weight-completed free outer image is separated.
In any fixed geometric weight \(m\), all odd
generators have weight at least three, so a
nonzero word has length at most \(m/3\).
Vanishing at every length therefore means
vanishing at every finite weight.
Together with \(\mathcal F_{n+1}=\ker\rho_{\leq n}\),
this proves the first equality in the theorem.  The second is
the Galois correspondence for the decreasing
closed normal subgroups \(\mathcal F_n\).

